% ===== 原文 pp.137–141（PDF p.147–151） =====
% [原文 p.137]
费米子场含量由形式
\begin{gather*}
  \rho,\tilde\rho\in\Omega^0(W,\mathit{ad}\,P)\\
  \lambda\in\Omega^1(W,\mathit{ad}\,P)\\
  \zeta\in\Omega^2(W,\mathit{ad}\,P)
\end{gather*}
组成。BRST 变分为
\begin{align}
  \delta_QA &= \lambda, \notag\\
  \delta_Q\varphi &= \ii\lambda, \notag\\
  \delta_Q\lambda &= 0, \notag\\
  \delta_Q\zeta &= -\ii\mathcal{F}, \tag{B.3}\\
  \delta_Q\rho &= 0, \notag\\
  \delta_Q\tilde\rho &= \ii P, \notag\\
  \delta_QP &= 0 \notag
\end{align}
其中 $\mathcal{F}$ 是复化联络 $\mathcal{A}=A+\ii\varphi$ 的曲率。我们用 $\dd_{\mathcal A}$ 记关于 $\mathcal{A}$ 的协变导数，用 $\dd_{\bar{\mathcal A}}$ 记关于 $\bar{\mathcal A}=A-\ii\varphi$ 的协变导数，并用 $\bar{\mathcal F}$ 记 $\bar{\mathcal A}$ 的曲率。微分算子 $\dd^\star_{\mathcal A}$ 是 $\star\dd_{\mathcal A}\star$，其中 $\star$ 是 3d Hodge 星。该理论的作用量在相差一个与度量无关的项的意义下是 BRST-正合的：
\begin{equation}
  \tilde S=-\frac{1}{2e^2}\delta_Q\int_W\Tr\left(\ii\zeta\wedge\star\bar{\mathcal F}
  +\ii\tilde\rho\wedge\star\left(P-2\dd^\star_{\mathcal A}\varphi\right)\right)
  +\int_W\Tr\left(\rho\wedge\dd_{\mathcal A}\zeta\right). \tag{B.4}
\end{equation}
该作用量具有一个鬼数对称性，在该对称性下 $\rho$ 与 $\lambda$ 的荷为 1，而 $\zeta$ 与 $\tilde\rho$ 的荷为 $-1$。$P$ 的运动方程读出 $P=\dd^\star_{\mathcal A}\varphi$；若把这个值代入 BRST 变换，则 BRST 算子只在模去费米子运动方程的意义下幂零。

把费米子 0-形式 $\rho$ 看作取值于 $\fg^*$ 而非 $\fg$ 更为自然，因为这样作用量中非 BRST-正合的部分就写成
\[
  \int_W\rho_a\wedge\dd_{\mathcal A}\zeta^a,
\]
它显然不依赖于 $\fg$ 上度量的选取。\footnote{对单李代数 $\fg$，可以借助 Killing 形式给出 $\fg$ 与 $\fg^*$ 的典范等同；但若 $\fg$ 含阿贝尔子代数，则度量并不唯一确定。}

% [Tutor Note: p.137 复化联络 $\mathcal{A}$ 是 BRST-不变的，这可由 (B.3) 一行算出：$\delta_Q\mathcal{A}=\delta_QA+\ii\delta_Q\varphi=\lambda+\ii\cdot\ii\lambda=0$。这正是 p.138 “对 BRST-不变联络 $\mathcal{A}$ 构造 Wilson 线”一句的依据；若把 (B.3) 中的 $\delta_Q\varphi=\ii\lambda$ 抄成 $\delta_Q\varphi=\lambda$，$\delta_Q\mathcal{A}$ 就不再为零，可据此校验该式。]

% [原文 p.138]
\subsection*{B-型 3d 规范理论的可观测量}

该拓扑规范理论中的局部可观测量是 $\rho$ 的规范不变函数，它们对应于外代数 $\Lambda^\bullet(\fg)$ 中关于 Ad 作用不变的元素。与 2d 情形不同，这里没有无序局部算子。

最简单的线算子可以构造为对 BRST-不变联络 $\mathcal{A}=A+\ii\varphi$ 的 Wilson 线。这类算子由 $G$ 的有限维表示标记。更一般地，我们可以考虑把生活在直线上的 1d TFT 与该 3d 规范理论耦合。该 1d TFT 的态空间是一个 $\Z$-分次向量空间 $V$。$V$ 的自同态也自然地带有分次。我们把该理论中生成 BRST 对称性的 1 次自同态记作 $T(\Phi)\in\End(V)$，把生成规范对称性的 0 次自同态记作 $R_a(\Phi)\in\End(V)$，其中 $\Phi$ 代表拓扑规范理论中的场。由于 $\rho$ 是该 3d 规范理论中唯一 BRST-不变的 0-形式，可以假定 $T$ 只是 $\rho$ 的函数。它还必须是幂零的：
\begin{equation}
  T(\Phi)=T(\rho), \tag{B.5}
\end{equation}
\begin{equation}
  T^2=0. \tag{B.6}
\end{equation}
由于规范对称性保持 $V$ 的分次，规范对称性的生成元 $R_a$ 可以假定与 $\rho$ 无关：
\begin{equation}
  R_a(\Phi)=R_a. \tag{B.7}
\end{equation}
由于该 3d 规范理论中规范对称性 $\delta_g$ 与 BRST 对称性 $\delta_Q$ 对易，$T$ 与 $R_a$ 必须满足如下关系：
\begin{align}
  0 &= [\delta_g(f),\delta] \notag\\
    &= [f,\rho]^a\frac{\partial T}{\partial\rho^a}+[f^aR_a,T] \tag{B.8}
\end{align}
其中 $f\in\fg$。为了构造与三元组 $(V,T,R)$ 相关联的线可观测量，我们对 $T$ 应用下降手续，从而得到以 $V$ 为纤维的分次向量丛上的一个联络 1-形式。按定义，后代联络 $\mathcal{N}$ 由方程
\[
  \delta\mathcal{N}=\dd T+[\mathcal{N},T]
\]
给出。利用关系 (B.8) 与费米子运动方程，我们得到
\begin{equation}
  \mathcal{N}=\frac{\ii}{2e^2}\star\bar{\mathcal F}^a\frac{\partial T}{\partial\rho^a}+A^aR_a. \tag{B.9}
\end{equation}

% [Tutor Note: p.138 此处说“$\rho$ 是唯一 BRST-不变的 0-形式”，严格说辅助场 $P$ 也满足 $\delta_QP=0$（(B.3) 末行）；但由 $P$ 的运动方程 $P=\dd^\star_{\mathcal A}\varphi$（p.137）可知 $P$ 并非独立自由度，故可取 $T$ 只依赖 $\rho$。这正是原文省略的一步。]

% [Tutor Note: p.138 (B.9) 的两项可分别解读：第二项 $A^aR_a$ 是一般情形下沿 Wilson 线的规范联络；当 1d TFT 取平凡（$T=1$，故 $\partial T/\partial\rho^a=0$）时 (B.9) 退化为 $\mathcal{N}=A^aR_a$，其超迹的和乐（p.139）正是表示 $R$ 中的 Wilson 线。第一项则是“下降”产生的部分，把 $\rho$ 换成了 $\star\bar{\mathcal F}^a$ 型的 1-形式。]

% [原文 p.139]
因此，$\mathcal{N}$ 沿 $W$ 中曲线 $\gamma$ 的和乐的超迹，是该拓扑规范理论中一个 BRST-不变且规范不变的圈算子。和乐本身则定义了一个线算子。

\subsection*{线算子的辫幺半范畴}

任何 3d TFT 中的线算子构成一个辫幺半范畴。上述线算子构成的子范畴是 DG-代数 $\Lambda^\bullet(\fg)$（微分为零）上 DG-模的 $G$-等变导出范畴。为看出这一点，考虑插入在两条 Wilson 线交接处的局部算子，这两条 Wilson 线对应三元组 $(V_1,T_1,R_1)$ 与 $(V_2,T_2,R_2)$。由于我们在寻找 BRST-不变的算子，可以假定它是取值于 $\Hom_\C(V_1,V_2)$ 的 $\rho$ 的函数 $\cO$，换言之是 $\Hom_\C(V_1,V_2\otimes\Lambda^\bullet(\fg))$ 的一个元素。BRST 算子作用在 $\cO$ 上为
\[
  \delta_Q\cO=T_2\cO\pm\cO T_1
\]
其中符号取正或取负取决于 $\cO$ 的总次数是奇还是偶。规范变换以显然的方式作用在 $\cO$ 上，并与 BRST 算子对易。线算子之间的态射空间是 $\delta_Q$ 在 $\Hom_\C(V_1,V_2\otimes\Lambda^\bullet(\fg))$ 的 $G$-不变部分上的上同调。

幺半结构在经典层面是显然的，由张量积给出。这个结果不可能有量子修正，因为规范耦合 $e^2$ 是一个无关参数。出于同样的理由，辫是平凡的。

还存在更一般的线算子。为看出这一点，我们可以使用维数约化的技巧，把 3d 理论中线算子的范畴等同于把 3d 理论在圆上紧化所得 2d 理论中膜的范畴。可以证明，约化给出的是以 $G_\C$ 为靶空间、与规范群为 $G$ 的 B-型规范理论耦合的 B-模型。从 3d 的角度看，$G_\C$ 参数化联络 $\mathcal{A}$ 沿紧化圆的和乐。规范群 $G$ 通过共轭作用在 $G_\C$ 上。正如 A.3 中所解释的，该 TFT 的膜范畴是凝聚层的等变导出范畴 $D^b_{G_\C}\big(\mathit{Coh}(G_\C)\big)$。上面讨论的线算子对应于该范畴中支撑在 $G_\C$ 单位元处的凝聚层。物理上，这来自这样的事实：对这类线算子，规范场 $\mathcal{A}$ 是非奇异的，因此绕该线算子的小圆上的和乐必须是平凡的。从数学的观点看，我们可以注意到：DG-代数 $\Lambda^\bullet(\fg)$ 的等变导出范畴通过 Koszul 对偶，等价于 DG-代数 $\mathrm{Sym}^\bullet(\fg^*)$ 的等变导出范畴（视其位于次数零）的一个由有限维 DG-模构成的全子范畴。后一范畴也可以看作 $G_\C$ 的等变导出范畴中“支撑”在 $G_\C$ 单位元处的全子范畴。也就是说，只关注 Wilson 型（即不是无序算子的）线算子，等价于只关注 $G_\C$ 上支撑在单位元处的等变层

% [Tutor Note: p.139 这里的对应关系可逐项核对：1d TFT 的态空间 $V$ 配上一个 1 次幂零算子 $T$（(B.6)）正是“DG-向量空间”的数据，$T^2=0$ 就是 DG-模的微分条件；而 (B.8) 则是 $T$ 与 $G$-作用 $R_a$ 的相容性（即微分与表示对易）。故“线算子范畴 = $\Lambda^\bullet(\fg)$ 上 DG-模的等变导出范畴”中的 $\Lambda^\bullet(\fg)$ 扮演了函数环的角色。]

% [Tutor Note: p.139 维数约化论证有一个简单的自洽检验：复化群 $G_\C$ 的复维数等于 $\dim G$，与“$G_\C$ 参数化和乐（$\mathcal{A}$ 的 1-形式沿圆的分量）给出 $\dim G$ 个额外标量自由度”一致；若误写成 $G$ 而非 $G_\C$，则维数只有一半，无法容纳 $\mathcal{A}=A+\ii\varphi$ 的实部与虚部两部分。]

% [原文 p.140]
元素。更一般地，人们还可以考虑 Gukov-Witten 型线算子，其 $\mathcal{A}$ 的和乐的共轭类是固定的；这类线算子可以看作 $D^b_{G_\C}\big(\mathit{Coh}(G_\C)\big)$ 中支撑在 $G_\C$ 的非平凡共轭类上的对象。

\section{带规范场的 Rozansky-Witten 模型}

我们通过把 Rozansky-Witten 模型与 B-型 3d 规范理论耦合，来对靶空间为 $X$ 的 Rozansky-Witten 模型做规范化。假设超 Kähler 靶空间 $X$ 容许一个相容的 $G$-作用。

设 $V_a$（$a=1,2,\cdots,\dim G$）是 $X$ 上对应 $G$-作用生成元的向量场。设 $\mu_+$、$\mu_-$ 与 $\mu_3$ 分别是相应于全纯辛形式 $\Omega$、反全纯辛形式 $\bar\Omega$ 与 Kähler 形式 $J$ 的矩映射，即
\begin{gather*}
  \dd\mu_{+a}=-\ii V_a(\Omega),\\
  \dd\mu_{-a}=-\ii V_a(\bar\Omega),\\
  \dd\mu_{3a}=\ii V_a(J).
\end{gather*}
场的 BRST 变分为
\begin{align}
  \delta_QA &= \lambda, & \delta_Q\phi^i &= 0, \notag\\
  \delta_Q\varphi &= \ii\lambda, & \delta_Q\phi^{\bar\imath} &= \eta^{\bar\imath}, \notag\\
  \delta_Q\lambda &= 0, & \delta_Q\eta^{\bar\imath} &= 0, \notag\\
  \delta_Q\zeta &= -\ii\mathcal{F}, & \delta_Q\chi^i &= \mathcal{D}\phi^i, \tag{B.10}\\
  \delta_Q\rho &= \ii\mu_+, \notag\\
  \delta_Q\tilde\rho &= \ii P, \notag\\
  \delta_QP &= 0 \notag
\end{align}
其中 $\mathcal{D}\phi^i=\dd\phi^i+\mathcal{A}^aV^i_a$ 是关于复化联络的规范协变导数。

带规范场 Rozansky-Witten 模型的作用量由
\[
  S=\int_W\left(\mathcal{L}_1+\mathcal{L}_2+\mathcal{L}_3+\mathcal{L}_4\right)
\]

% [Tutor Note: p.140 (B.10) 可视为 (B.3) 与 RW 模型 BRST 变分（p.136）的“并集＋协变化”：$\delta_Q\chi^i$ 由 RW 的 $\dd\phi^i$ 变成 $\mathcal{D}\phi^i$，而 $\delta_Q(\text{规范扇区})$ 与 (B.3) 逐行一致（如 $\delta_Q\varphi=\ii\lambda$、$\delta_Q\zeta=-\ii\mathcal{F}$）。据此可校验 (B.10) 没有漏项。]

% [原文 p.141]
给出，其中
\begin{align}
  \mathcal{L}_1 &= -\frac{1}{2e^2}\delta_Q\Tr\left(\ii\zeta\wedge\star\bar{\mathcal F}
    +\ii\tilde\rho\wedge\star\left(P-2\dd^\star_{\mathcal A}\varphi-2e^2\mu_3\right)\right), \tag{B.11}\\
  \mathcal{L}_2 &= \delta_Q\left(g_{i\bar\jmath}\chi^i\wedge\star\mathcal{D}\phi^{\bar\jmath}\right), \tag{B.12}\\
  \mathcal{L}_3 &= \delta_Q\left(\frac{\ii}{2}e^2\star\rho^a\mu_{-a}\right), \tag{B.13}\\
  \mathcal{L}_4 &= \rho_a\wedge\dd_{\mathcal A}\zeta^a+\ii\Omega_{ij}\chi^i\wedge\zeta^aV^j_a
    +\frac12\Omega_{ij}\chi^i\wedge\mathcal{D}\chi^j \tag{B.14}\\
  &\quad+\frac16\Omega_{ij}\mathcal{R}^j{}_{kl\bar m}\chi^i\wedge\chi^k\wedge\chi^i\wedge\eta^{\bar m} \notag
\end{align}
其中 $\mathcal{D}\chi^i=\dd\chi^i+\Gamma^i_{jk}\dd\phi^j\wedge\chi^k+\mathcal{A}^a\partial_jV^i_a\wedge\chi^j+\mathcal{A}^a\Gamma^i_{jk}V^k_a\wedge\chi^j$。

带规范场 Rozansky-Witten 模型的局部可观测量是 $\rho_a$、$\phi^i$、$\phi^{\bar\imath}$ 与 $\eta^{\bar\imath}$ 的 BRST 不变且规范不变的函数，它们对应于 $\Lambda^\bullet(\fg)\otimes\Omega^{0,\bullet}(X)$ 关于如下幂零算子的上同调 [11]：
\[
  \delta=\ii\mu_{+a}T^a+(-1)^l\dbar_X
\]
其中 $T^a$ 是 $\fg$ 的一组基的元素，$l$ 是 $\Lambda^\bullet(\fg)$ 中元素的次数。

% [Tutor Note: p.141 (B.14) 末项的 $\chi$ 指标与 p.136 有同一处排印错误：印作 $\chi^i\wedge\chi^k\wedge\chi^i$（$i$ 重复），按形式次数（3-流形上应为 3-形式）与 $\mathcal{R}^j{}_{kl\bar m}$ 的指标缩并，应为 $\chi^i\wedge\chi^k\wedge\chi^l\wedge\eta^{\bar m}$。与 p.136 的对应项（未规范化版本）比对即可确认两处应一致。译文照原文排印。]

% [Tutor Note: p.141 末式 $\delta=\ii\mu_{+a}T^a+(-1)^l\dbar_X$ 的幂零性需要两条事实：(i) $\dbar_X^2=0$ 且 $\mu_+$ 全纯（故 $\dbar_X\mu_+=0$）；(ii) $T^a$ 给出 $\fg$ 在 $\Lambda^\bullet(\fg)$ 上的表示，且 $\ii\mu_{+a}T^a$ 的平方项与 $\dbar_X$ 的交叉项相消（等价于 $G$-作用与 BRST 的相容性）。原文只断言其幂零，未给证明；核对时可与 (B.10) 的 $\delta_Q^2=0$ 结构对照。]
